\documentclass{article}
\usepackage{../math_template}

\author{Jonathan Kasongo}
\title{Imperial-Cambridge Math Contest 2022 --- P1/6}

\begin{document}
\maketitle

\begin{problem}{Imperial-Cambridge Math Contest 2022 --- P1/6}
Two straight lines divide a square of side length 1 into four regions. Show
that at least one of the regions has a perimeter greater than or equal to 2.
\end{problem}

\begin{solution}{Solution}
Let $Q_1 Q_2$ and $P_1 P_2$ divide the square $ABCD$ into four regions,
and let $\gamma$ be the intersection of $P_1 P_2$ and $Q_1 Q_2$.


\begin{center}
\begin{asy}
/* Geogebra to Asymptote conversion, documentation at artofproblemsolving.com/Wiki go to User:Azjps/geogebra */
import graph; size(100);
real labelscalefactor = 0.5; /* changes label-to-point distance */
pen dps = linewidth(0.7) + fontsize(10); defaultpen(dps); /* default pen style */
pen dotstyle = black; /* point style */
real xmin = -3.365825512563786, xmax = 2.912545802957398, ymin = -2.5070460973012123, ymax = 1.870713201303435;  /* image dimensions */
pen xdxdff = rgb(0.49019607843137253,0.49019607843137253,1.); pen ffqqtt = rgb(1.,0.,0.2);
pair A = (-1.,1.), B = (1.,1.), C = (1.,-1.), D = (-1.,-1.), P_1 = (-1.,-0.5), P_2 = (1.,0.5), Q_1 = (-0.5,1.), Q_2 = (0.1275873205861335,-1.), gamma  = (-0.16095328023618855,-0.08047664011809427);

filldraw(A--B--C--D--cycle, xdxdff + opacity(0.10000000149011612),  xdxdff);
/* draw figures */
draw(A--B,  xdxdff);
draw(B--C,  xdxdff);
draw(C--D,  xdxdff);
draw(D--A,  xdxdff);
draw(P_1--P_2,  linetype("4 4") + ffqqtt);
draw(Q_1--Q_2,  linetype("4 4") + ffqqtt);
/* dots and labels */
dot(A,dotstyle);
label("$A$", (-0.9815101802880369,1.0498833328150634), NE * labelscalefactor);
dot(B,dotstyle);
label("$B$", (1.0217055701895397,1.0498833328150634), NE * labelscalefactor);
dot(C,dotstyle);
label("$C$", (1.0217055701895397,-0.9533324176625092), NE * labelscalefactor);
dot(D,dotstyle);
label("$D$", (-0.9815101802880369,-0.9533324176625092), NE * labelscalefactor);
dot(P_1,dotstyle);
label("$P_1$", (-0.9815101802880369,-0.45008553400594825), NE * labelscalefactor);
dot(P_2,dotstyle);
label("$P_2$", (1.0217055701895397,0.5466364491585025), NE * labelscalefactor);
dot(Q_1,dotstyle);
label("$Q_1$", (-0.47826329663147504,1.0498833328150634), NE * labelscalefactor);
dot(Q_2,dotstyle);
label("$Q_2$", (0.14713088888347572,-0.9533324176625092), NE * labelscalefactor);
dot(gamma ,dotstyle);
label("$\gamma$", (-0.1411367435023219,-0.039670599761762654), NE * labelscalefactor);
clip((xmin,ymin)--(xmin,ymax)--(xmax,ymax)--(xmax,ymin)--cycle);
/* end of picture */
\end{asy}
\end{center}

If $P_1 P_2$ and $Q_1 Q_2$ both intersect opposite sides of the square
then they both will have length $ \geq 1$. Now the total perimeter $P$ of
the four regions will be
$$
P = 4 + 2(Q_1 \gamma + Q_2 \gamma + P_1 \gamma + P_2 \gamma)
\geq 4 + 2(1+1) = 8.
$$
Now the region with the largest perimeter must have perimeter $\geq 2$,
because otherwise the total perimeter would be $P < 2+2+2+2 = 8$, a
contradiction.\vspace{0.2cm}

Now if exactly one of these lines intersect opposite sides of the square,
say WLOG $P_1 P_2$. Now WLOG let $Q_1 Q_2$ be such that $Q_1$ lies on the
opposite side of whatever side $P_2$ lies on.
Then one of the four regions is bounded by
$P_2 \gamma$, $Q_1 \gamma$ and one of the sides of the square.
Notice that $P_2 \gamma + Q_1 \gamma \geq P_2 Q_1 \geq 1$.
The perimeter
of this region is $\geq 1+1 = 2$. \vspace{0.2cm}

Now finally if none of the lines intersect opposite sides of the square
then one of the regions includes two adjacent sides of the square, which
means the perimeter of that region is at least $1+1=2$. $\blacksquare$
\end{solution}

\end{document}
